% Originally: content/week-11.tex, \subsection{Sneak peek at Stokes' theorem} --- promoted to a section

\section{Sneak peek at Stokes' theorem}
\label{sec:stokes_sneak_peek}

\begin{ainote}
Before the formal introduction of Stokes' theorem (\cref{sec:stokes_general}), Sascha presented introductory examples to build intuition. The classical version of Stokes' theorem states that for a surface $A \subset \mathbb{R}^3$ and a vector field $F$,
\[\int_A (\curl F)\cdot \nu \mathop{d\text{vol}}_2 = \int_{\partial A} F \cdot \tau \,dL,\]
where $\nu$ is the unit normal to $A$ and $\tau$ is the unit tangent vector to the boundary curve $\partial A$.
\end{ainote}

\begin{example}[Stokes' theorem on a hemisphere]
\label{ex:stokes_hemisphere}
Let $A := \{(x,y,z) \in \mathbb{R}^3 \mid x^2+y^2+z^2=1,\ z \geq 0\}$ be the upper hemisphere.
Given the vector field $F(x,y,z) := (z^2, x^2, y^2)\transp$, we can evaluate both sides of Stokes' theorem.
\begin{enumerate}[label=\textbf{\arabic*.}]
    \item \textbf{The rotation of $F$:} $\curl F = \nabla \times F = (2y, 2z, 2x)\transp$.
    \item \textbf{The surface normal:} On the unit sphere, the outward normal is simply $\nu(x,y,z) = (x,y,z)\transp$.
    \item \textbf{The boundary $\partial A$:} The boundary of the upper hemisphere is the unit circle in the $xy$-plane, $x^2+y^2=1$ with $z=0$. This can be parametrized by the closed path $\gamma : [0,2\pi] \to \mathbb{R}^3$,
    \[\gamma(t) = \begin{pmatrix} \cos t \\ \sin t \\ 0 \end{pmatrix}.\]
\end{enumerate}
Evaluating the path integral (the right-hand side of Stokes' theorem):
\[\int_{\partial A} F \cdot \tau \,dL = \int_0^{2\pi} F(\gamma(t)) \cdot \gamma'(t) \,dt.\]
Since $z=0$ on $\gamma$, we have $F(\gamma(t)) = (0, \cos^2 t, \sin^2 t)\transp$. The derivative is $\gamma'(t) = (-\sin t, \cos t, 0)\transp$.
\[\int_0^{2\pi} \begin{pmatrix} 0 \\ \cos^2 t \\ \sin^2 t \end{pmatrix} \cdot \begin{pmatrix} -\sin t \\ \cos t \\ 0 \end{pmatrix} dt = \int_0^{2\pi} \cos^3 t \,dt = 0.\]
\end{example}

\begin{example}[Path integral using Stokes' theorem]
\label{ex:stokes_path_integral}
Compute $\oint_\Gamma F \cdot \tau \,dL$ where $F(x,y,z) := (x, y, z)\transp$ and $\Gamma$ is the unit circle in the $xy$-plane (oriented counter-clockwise).

Instead of computing the line integral directly, we can use Stokes' theorem. The curve $\Gamma$ bounds the unit disk $A = \{(x,y,0) \mid x^2+y^2 \leq 1\}$.
First, we compute the rotation of $F$:
\[\curl F = \nabla \times F = \begin{pmatrix} \partial_y(z) - \partial_z(y) \\ \partial_z(x) - \partial_x(z) \\ \partial_x(y) - \partial_y(x) \end{pmatrix} = \begin{pmatrix} 0 \\ 0 \\ 0 \end{pmatrix}.\]
By Stokes' theorem, the line integral along $\Gamma$ bounds a surface $A$ on the sphere (for example, the upper hemisphere). The surface integral is:
\[\oint_\Gamma F \cdot \tau \,dL = \int_A (\curl F) \cdot \nu \mathop{d\text{vol}}_2 = 0.\]
\end{example}

% Supplement: Sascha Brack/Class Notes/Week_11_Notes_Monday.pdf, pp. 5--10
